
%(BEGIN_QUESTION)
% Copyright 2006, Tony R. Kuphaldt, released under the Creative Commons Attribution License (v 1.0)
% This means you may do almost anything with this work of mine, so long as you give me proper credit

Hvor mange elektroner, protoner og nøytroner inneholder ett enkelt atom av karbon-12 ($^{12}$C)? (Anta elektrisk nøytralt atom.)

\vskip 10pt

Hvor mange elektroner, protoner og nøytroner inneholder et gjennomsnittlig gullatom (Au)? (Anta elektrisk nøytralt atom.)

\underbar{file i00553}
%(END_QUESTION)





%(BEGIN_ANSWER)

For karbon-12 representerer ``12'' atomets {\it atommasse} (summen av protoner og nøytroner i kjernen). Siden {\it atomnummeret} for karbon er 6, vet vi at kjernen har 6 protoner. Da gjenstår 6 nøytroner (12 minus 6). For at atomet skal være elektrisk nøytralt, må det være 6 elektroner som balanserer de 6 protonene.

\vskip 10pt

For gull representerer den gjennomsnittlige {\it atommassen} 197 summen av protoner og nøytroner i kjernen. Siden {\it atomnummeret} for gull er 79, vet vi at kjernen har 79 protoner. Da gjenstår 118 nøytroner (197 minus 79). For at atomet skal være elektrisk nøytralt, må det være 79 elektroner som balanserer de 79 protonene.

%(END_ANSWER)





%(BEGIN_NOTES)


%INDEX% Kjemi, grunnprinsipper: proton, nøytron, elektron

%(END_NOTES)

